\centerline{\bf \Huge Сравнения по модулю}

\begin{flushright}
        \scriptsize{}
\end{flushright}

В этом файле мы постараемся досконально изучить такую важную в теории чисел тему, как сравнения по модулю. По сути, это более удобная форма для записи остатков. Собственно, начнем с определения и простейших свойств.

\textbf{Определение.} Целые числа, разность которых делится на $m$, называются сравнимыми по модулю $m$. Другими словами, если $a-b$ делится на $m$, то $a$ и $b$ сравнимы по модулю $m$. Запись: $a \equiv b\ (mod\ m)$.

Из этого определения следует, что если $a \equiv b\ (mod\ m)$, то $a$ и $b$ дают одинаковые остатки при делении на $m$. И ради этого уже стоит вводить сравнения: до этого удобного способа написать, что у чисел равные остатки, не было!

Конечно, сравнения были бы бесполезны, если у записи выше не было каких-то свойств. И они есть, более того, эти свойства напоминают нам свойства \textbf{обычных равенств}. Поэтому их легко запомнить и применять.

Очень важно, что эти свойства нам не даны свыше, а мы можем их доказать из определения сравнений, чем и займемся сразу после формулировки всех свойств.

\textbf{Свойства сравнений.}

\begin{itemize}
    \item Сравнения можно умножать на число:
    
    если $a \equiv b\pmod{m}$, то $k a \equiv k b\pmod{m}$.
    \item Сравнения можно складывать и вычитать:
    
    если $a \equiv b\pmod{m}$ и $c \equiv d\ (mod\ m)$, то $a \pm c \equiv b \pm d\pmod{m}$.
    \item Сравнения можно перемножать:
    
    если $a \equiv b\pmod{m}$ и $c \equiv d\pmod{m}$, то $a c \equiv b d\pmod{m}$.
    \item Сравнения можно возводить в степень:
    
    если $a \equiv b\pmod{m}$, то $a^k \equiv b^k\pmod{m}$.
\end{itemize}


\newpage

\hspace{3mm}

\textbf{Доказательство свойств:}

1. Из того, что $a \equiv b \pmod{m}$ следует, что разность $a-b: m$. Умножим разность на число $k$, тогда делимость не пропадет: $k(a-b) \vdots m$. Раскроем скобки: $k a-k b: m$. Теперь от делимости вернемся к сравнениям: из последней делимости по определению сравнений следует, что $k a \equiv k b \pmod{m}$, что и требовалось доказать.

2. Из определения сравнение по модулю следует, что $a\pm b \ \vdots \ m$ и $c\pm d\ \vdots \ m$. Тогда

\[
(a\pm c)-(b\pm d) \ \vdots \ m
\]

откуда и следует, что $a \pm c \equiv b \pm d\pmod{m}$.

3. По первому свойству, которое мы уже доказали, первое данное в условии сравнение можно умножить на число $c$, тогда мы получим $a c \equiv b c\pmod{m}$. Аналогично второе сравнение можно умножить на $b: b c \equiv b d\pmod{m}$. Объединяя выведенные условия, имеем $a c \equiv b c \equiv b d\pmod{m}$, или, отбрасывая промежуточное сравнение, $a c \equiv b d\pmod{m}$, что и требовалось доказать.

4. Последнее свойство следует из третьего, когда мы перемножаем сравнение само с собой $k$ раз. чтобы получить $k$-е степени.

Итак, теперь, когда мы доказали все свойства, пора сформулировать главный

\textbf{Вывод.}
В сравнениях мы можем заменять число (\textbf{не степень!!}) на любое число, дающее тот же остаток по рассматриваемому модулю.

\textbf{Пример.} Докажите, что при любом натуральном $n$ выполнено $2^{4 n}-1: 15$.
< Перепишем условие в виде сравнения. Нам нужно доказать, что $2^{4 n} \equiv 1(\bmod 15)$. Распишем левую часть по-другому, возведя двойку в 4-ю степень:

$$
2^{4 n} \equiv\left(2^4\right)^n \equiv 16^n \equiv 1^n \equiv 1 \quad(\bmod 15)
$$


В последних двух сравнениях мы заменили 16 на 1 по модулю 15 , потому что они сравнимы, а дальше заметили, что 1 в любой степени --- по прежнему 1 . В итоге сравнение доказано.

\newpage

\hspace{3mm}


\centerline{\bf \Huge Лайфхаки при решении задач}

В первую очередь мы должны с вами раз и навсегда запомнить свойства степеней:

$$
\begin{aligned}
& a^m \cdot a^n=a^{m+n} \\
& a^m \div a^n=\frac{a^m}{a^n}=a^{m-n} \\
& \left(a^m\right)^n=a^{m \cdot n} \\
& (a \cdot b)^n=a^n \cdot b^n \\
& \left(\frac{a}{b}\right)^n=\frac{a^n}{b^n} \\
& 1^n=1 \\
& a^0=1
\end{aligned}
$$

Далее мы опишем главные приёмы, которые стоит использовать при решении задач на сравнения по модулю.

\begin{enumerate}
    \item \textbf{Замена числа на его остаток.}
    
    \textit{Пример.} Докажите, что число $1000 \cdot 1001 \cdot 1002 \cdot 1003$ - 24 делится на 999.

    \textbf{Решение.} Заменим каждое число на его остаток. 

    $$
    1000 \equiv 1 \pmod{999} \quad 1001 \equiv 2 \pmod{999}
    $$

    $$
    1002 \equiv 3 \pmod{999} \quad 1003 \equiv 4 \pmod{999}
    $$
    
    Перемножим эти четыре сравнения, благо мы можем это сделать благодаря нашим свойствам и получим, что 

    $$
    1000\cdot10001\cdot1002\cdot1003 \equiv 1\cdot2\cdot3\cdot4 = 24 \pmod{999}
    $$
    
    что и требовалось доказать.

    \item \textbf{Замена числа на отрицательный остаток.}
    Иногда бывает удобным рассматривать не обычную систему остатков а так называемую отрицательную систему.

    Например, не очень удобно рассматривать число 1003 по модулю 1004. Но зато благодаря сравнениям мы понимаем, что 

    $$
    1003 \equiv -1 \pmod{1004} \text{, т.к. } 1003 - (-1) = 1004 \equiv 0 \pmod{1004}
    $$

    Иначе говоря, остатки $1, 2, 3, \ldots, n-2, n-1 $ можно заменять на $-(n-1), -(n-2), -(n-3), \ldots, -2, -1 $ по модулю $n$.
    
    \textit{Пример 1.}
    Докажите, что число $1000 \cdot 1001 \cdot 1002 \cdot 1003-24$ делится на 1004.

    \textbf{Решение.}
    Заменим каждое число на его отрицательный остаток по модулю $1004$.


    $$
    1000 \equiv -1 \pmod{1004} \quad 1001 \equiv -2 \pmod{1004}
    $$

    $$
    1002 \equiv -3 \pmod{1004} \quad 1003 \equiv -4 \pmod{1004}
    $$

    Опять перемножим все сравнения и получим требуемое.
    
    \textit{Пример 2.} Докажите, что число $5^{70}+6^{70}$ делится на 61 .

    \textbf{Решение.}
    Заметим, что $5^2 = 25 \equiv -36 \equiv -6^2 \pmod{61}$

    Тогда по свойствам степеней и сравнений по модулю получаем, что

    $$
    5^{70} = (5^2)^{35} \equiv (-6^2)^{35} = -6^{70} \pmod{61} 
    $$

    а значит исходное выражение даёт остаток 0 при делении на $61$.

    \item \textbf{Возводить единичку в степень.}

    Зачастую очень полезно возводить в степень сравнения вида $a \equiv 1 \pmod{m}$, так как для любого $k$ будет выполнено, что $a^k \equiv 1 \pmod{m}$. 

    \textit{Пример 1.} С помощью этого можно с лёгкость доказать признак равноостаточности по модулю 9 (он же всем известный признак делимости)

    Число $\overline{a_1a_2\ldots a_n}$ (то есть $n$-значное число состоящее из цифр $a_1,a_2, \ldots, a_n$) представляется в виде

    $$
    \overline{a_1a_2\ldots a_n} = a_1\cdot 10^{n-1} + a_2 \cdot 10^{n-2} + \ldots a_n \cdot10^0
    $$

    Давайте проверим какой достаток это число при делении на 9. Заметим, что $10\equiv1 \pmod{9}$, а значит и любая степень десятки даёт остаток 1. Стало быть 

    $$
    \overline{a_1a_2\ldots a_n} = a_1\cdot 10^{n-1} + a_2 \cdot 10^{n-2} + \ldots a_n \cdot10^0 \equiv a_1 + a_2 +\ldots + a_n \pmod{9}
    $$

    то есть число даёт такой же остаток при делении на 9, как и его сумма цифр. 

    \newpage

    \hspace{3mm}

    \textit{Пример 2.}
    Некоторые умные школьники знают признак делимости на 11. Звучит он так: если сумма цифр на четных минус сумма цифр на нечётных местах делится на 11, то и само число делится на 11. Знать то можно много чего, а вот как доказывается этот факт? Возьмём разложение из предыдущего пункта и опять попробуем разобраться, какой остаток дают степени десятки. Само число 10 даёт, собственно, остаток 10 при делении на 11. Но мы с вами уже научились трюку отрицательных остатков, поэтому $10 \equiv -1 \pmod{11}$ и тогда $10^k \equiv (-1)^k \pmod{11}$. Значит в чётных степенях 10 даёт остаток 1, а в нечётных даёт <<остаток>> -1. То есть 

    $$
    \overline{a_1a_2\ldots a_n} = a_1\cdot 10^{n-1} + a_2 \cdot 10^{n-2} + \ldots a_n \cdot10^0 \equiv -a_1 + a_2 -a_3\ldots + (-1)^n\cdot a_n \pmod{11}
    $$

    что и требовалось доказать.
    
    \textit{Пример 3.}
    Найдите остаток от деления $47^{101}$ на 31

    \textbf{Решение.}
    Сразу же возводить 47 в степень 101 довольно трудоёмко. Давайте посмотрим на маленькие степени.

    $$
    47 \equiv 16 \pmod{31}
    $$

    $$
    47^2 \equiv 16^2 = 256 \equiv 8 \pmod{31}
    $$

    $$
    47^3 \equiv 47^2\cdot47 = 16 \cdot 8 = 128 \equiv 4 \pmod{31}
    $$

    $$
    47^4 \equiv (47^2)^2 \equiv 8^2  = 64 \equiv 2\pmod{31}
    $$

    $$
    47^5 \equiv 47^2 \cdot 47^3 \equiv 8\cdot4 = 32 \equiv 1\pmod{31}
    $$

    Не прошло и три годы, как мы дошли до единички. На самом-то деле мы бы дошли до неё в любом случае, но это уже более серьёзный факт, который нам предстоит пройти, когда мы повзрослеем. Теперь возведём последнее сравнение в удобную для нас степень, а именно 20-ю. Получим, что $47^{100} = (47^{5})^{20} \equiv 1 \pmod{31}$ и теперь это очень сильно поможет нам найти искомый остаток

    $$
    47^{101} \equiv 47^{100}\cdot47 \equiv 1\cdot47 \equiv16 \pmod{31}
    $$
\end{enumerate}
